% IV. Ръководство за потребителя
% Ориентировъчен обем: 9–10 стр.
%
% Пише се от гледна точка на ЧОВЕКА, който използва системата, не на този, който я
% поддържа. Затова тук няма имена на файлове, класове и входни точки: те са в
% Раздел III. Всяко твърдение тук е проверимо чрез самия интерфейс.
\section{Ръководство за потребителя}
\label{sec:userguide}

\subsection{Достъп до системата и общ изглед}
\label{subsec:ug-access}

Системата се отваря в браузър на адреса, на който е публикувана, и не изисква
инсталиране или добавка. Регистрация е необходима само за запазването на проекти и за
запитванията към инженерния отдел; \emph{самото изчисление е достъпно без вход в
системата}.

Всяка страница има еднаква заглавна лента с четири раздела, а вдясно бутоните за
профил. Разделите са: \textbf{Load Calculator} (изчисляване на товарите),
\textbf{My Projects} (запазените изчисления), \textbf{Technical Documentation}
(приложното ръководство и типовите детайли за закрепване) и \textbf{Engineering
Support} (запитване към инженерния отдел). На екран с ширина под 960 px разделите се
подреждат в две колони под знака на марката; съдържанието и последователността на
действията остават същите (§\ref{subsec:ug-mobile}).

Основната страница е калкулаторът \figref{fig:ss-calculator}: изборите стоят в
неподвижен панел вляво, а резултатът заема останалата ширина.

\screenshot{calculator}{Страница „Load Calculator“: избори вляво, таблица с товарите и схема на съоръжението вдясно}

\subsection{Изчисляване на товарите}
\label{subsec:ug-calc}

Изчислението се задава с шест избора в панела \textbf{Your inputs} вляво. Редът е
съществен: всеки следващ избор се ограничава от предходните, така че системата
предлага само комбинации, за които съществуват таблични данни.

\begin{enumerate}[topsep=2pt,itemsep=3pt,leftmargin=*]
\item \textbf{Units.} Мерна система: \emph{Metric (kN\,$\cdot$\,m)} или
      \emph{Imperial (lb\,$\cdot$\,ft)}. Изборът сменя както единиците, така и
      нормативната система, по която се определят частните коефициенти.
\item \textbf{Structure type.} \emph{Climbing wall} (катерачна стена с осигурителни
      точки) или \emph{Boulder wall} (боулдер стена без осигурителни точки). Под избора
      се изписва приложимата норма, например „EN 12572-1 (ACS with protection points)“.
\item \textbf{Climbing-surface height.} Височина на катерачната повърхност. За
      катерачна стена се предлагат стойности от 8 до 16 m, за боулдер стена 4 и 5 m.
\item \textbf{Attachment scheme.} Брой нива на закрепване. Предлагат се само схемите,
      които съществуват за избраната височина: например при 12 m това са 2 и 3 нива, а
      при 9 m само 2. За боулдер стена изборът не се показва, защото точката е една.
\item \textbf{Column span $\cdot$ A.} Разстояние между колоните на съществуващата
      сграда, от 4 до 8 m.
\item \textbf{Overhang $\cdot$ X.} Наклон на стената, тоест разликата между горния и
      долния контур, от 1 до 7 m.
\end{enumerate}

За катерачна стена се показва и седми избор, \textbf{Live load force level}. Той
определя за кое ниво на закрепване се взема максималният полезен товар; останалите нива
се отчитат едновременно с него. Изборът се прави отделно, защото не съществува едно
състояние, в което всички нива са едновременно максимално натоварени.

Под изборите стои превключвателят \textbf{Show factored design values}. При изключено
състояние (по подразбиране) се показват характеристичните стойности от таблиците. При
включено те се умножават с частните коефициенти на избраната нормативна система,
изписани до самия превключвател, например „$\times$1.35 DL, $\times$1.5 LL“.

Бутонът \textbf{Reset} в горния десен ъгъл на панела връща всички избори към началното
им състояние и изчиства въведените стойности за носимост.

\textbf{Ако вместо резултат се появи покана за избор.} Резултатът се показва едва
когато всички величини са избрани съзнателно. След смяна на типа или на височината
някои от следващите избори може да престанат да са валидни; системата ги привежда към
най-близката допустима стойност и ги маркира като неизбрани, тоест иска потвърждение.
Достатъчно е съответната стойност да бъде натисната отново.

\subsection{Разчитане на резултата}
\label{subsec:ug-read}

Заглавието на резултата повтаря избора с думи, например „Climbing wall 12 m“, а под
него, с по-дребен шрифт, стоят броят нива, разстоянието между колоните, наклонът и
мерната единица, в която са изразени числата. Този ред е контролен: ако той не описва
търсения случай, таблицата под него също не го описва.

Таблицата има три колони. Първата назовава точката, втората дава постоянния товар
(\textbf{DL}, черно), третата полезния (\textbf{LL}, червено). Означенията на точките
са същите като в еталонните таблици и са обяснени в \tabref{tab:symbols}. За бърза
справка:

\begin{itemize}[topsep=2pt,itemsep=1pt,leftmargin=*]
\item \code{RZ0} и \code{RX0} са реакциите в основата, вертикална и хоризонтална;
\item \code{RX1}, \code{RX2}, \code{RX3} са хоризонталните реакции на нивата на
      закрепване, с изписана височина на всяко ниво;
\item \code{LX1}, \code{LX2}, \code{LX3} са товарите върху \emph{колона на сградата} и
      вече отчитат разстоянието между колоните.
\end{itemize}

\textbf{За преценка дали сградата издържа, се гледат стойностите \code{LX}}, не
\code{RX}. Величините с представка \code{R} описват какво се случва в самото
съоръжение.

Отрицателният знак не означава грешка: той е посока по координатната система, а не
по-малка големина. Тире вместо число означава, че за тази точка таблицата не дава
стойност, а не че стойността е нула.

Вдясно от таблицата стои схема на съоръжението, изчертана от текущия избор. Тя носи
размерните линии, точките на закрепване и стрелките на товарите с техните стойности.
Схемата се увеличава с бутоните $+$ и $-$ в горния ѝ десен ъгъл, връща се към
първоначалния мащаб с кръговата стрелка и се придвижва с влачене. При тесен екран
схемата застава под таблицата и запазва възможността за увеличение, вместо да бъде
смалена до нечетливост.

\subsection{Двете схеми на закрепване}
\label{subsec:ug-options}

Над резултата стоят два раздела, които представят два различни начина за закрепване на
съоръжението към сградата. Те не са варианти на едно и също изчисление, а два различни
инженерни подхода, и дават различни таблици:

\begin{itemize}[topsep=2pt,itemsep=2pt,leftmargin=*]
\item \textbf{Option 1: Single-point attachment.} Съоръжението се закрепва пряко към
      съществуващата конструкция в отделни точки. Показват се реакциите в основата и на
      всяко ниво на закрепване.
\item \textbf{Option 2: Walltopia support beams.} Между колоните на сградата се
      поставят греди на Walltopia и съоръжението се закрепва към тях. Показват се
      товарите върху колоните на сградата, \code{LX1} до \code{LX3}.
\end{itemize}

Под таблицата на всяка от двете схеми се избират условията на закрепване: за първата
схема вида на подовата плоча (\emph{Hollow panel slab} или \emph{Solid concrete slab})
и метода на закрепване, за втората вида на носещата колона (\emph{Solid concrete
column} или \emph{Steel column}). Изборът отваря списък с типови детайли; всеки от тях
се разглежда с бутона \textbf{View}, а пълният чертеж с \textbf{View full detail}.
Точно \emph{един} детайл от списъка се потвърждава с бутона \textbf{Select}, който след
натискане показва \textbf{Selected}. До този избор схемата на съоръжението показва само
очертанието; след него точките на закрепване в схемата стават посочваеми и показват
стойността на товара при преглед.

\subsection{Проверка срещу носимостта на съществуващата конструкция}
\label{subsec:ug-capacity}

Разделът \textbf{Check capacity} под таблицата приема носимостта, която съществуващата
конструкция може да поеме, и я сравнява с получения товар. Полетата зависят от избраната
схема: при пряко закрепване се въвеждат носимостите в основата, а при схемата с греди
по едно поле за всяка колона (\code{LX1}, \code{LX2}, \code{LX3}).

Стойностите се въвеждат в мерната единица, изписана вдясно от полето, и се потвърждават
с бутона \textbf{Check capacity}. За всяка проверена точка системата извежда изискваната
стойност, въведената носимост и оценка: \emph{Applicable}, когато товарът се побира, и
\emph{Exceeds capacity}, когато го надвишава. Когато са попълнени всички полета, се
показва и обща оценка.

Оценката е сравнение на две числа и не е конструктивна проверка. Тя не отчита
комбинации от въздействия, състоянието на съществуващата конструкция и коефициентите за
конкретния случай. Резултатът \emph{Applicable} означава, че предварителните товари се
побират в заявената носимост, не че закрепването е допустимо.

\subsection{Профил: регистрация и вход}
\label{subsec:ug-account}

Бутонът \textbf{Register} в заглавната лента отваря прозорец с полета за име,
електронна поща и парола. Паролата трябва да е поне 8 знака. Един адрес може да има
само един профил; при повторна регистрация системата съобщава, че профил с този адрес
вече съществува. Бутонът \textbf{Log in} отваря същия прозорец в режим на вход, а
\textbf{Log out} прекратява сесията.

Сесията се пази седем дни и остава при затваряне на браузъра. При съобщение „Not
authenticated“ по време на работа сесията е изтекла и е достатъчно да се влезе отново;
незаписаното изчисление се запазва в браузъра и остава на екрана.

\subsection{Проекти: запазване, отваряне и промяна}
\label{subsec:ug-projects}

Бутонът \textbf{Save as project} под резултата запазва текущото изчисление. Отваря се
кратка форма с три полета:

\begin{itemize}[topsep=2pt,itemsep=1pt,leftmargin=*]
\item \textbf{Project name.} Задължително. Предлага се име по подразбиране, съставено от
      избора, например „Climbing wall 12 m $\cdot$ 3-level“.
\item \textbf{Tags.} Етикети, разделени със запетая, до шест на брой. Използват се за
      филтриране в списъка, например \emph{gym}, \emph{EU}, \emph{zone-2}.
\item \textbf{Custom properties.} Произволни двойки „свойство: стойност“, например
      „Client: Acme Climbing Ltd“ или „Ref: WT-2026-001“. Нов ред се добавя с
      \textbf{+ Add property}.
\end{itemize}

Отворен за редакция проект се обозначава с етикет \textbf{Editing} до бутоните. От този
момент \textbf{Save changes} записва промените в същия проект, а \textbf{Save as new}
създава копие и оставя оригинала непроменен. Стрелката до етикета излиза от режим на
редакция, като запазва числата на екрана като ново, незаписано изчисление.

\subsection{Списък с проекти}
\label{subsec:ug-dashboard}

Разделът \textbf{My Projects} показва запазените изчисления като карти
\figref{fig:ss-dashboard}. Всяка карта носи името, датата на последна промяна,
обобщение с водещия товар и оценката от проверката за носимост, етикетите и
потребителските свойства.

\screenshot{dashboard}{Страница „My Projects“: запазените изчисления като карти с обобщение, етикети и действия}

Действията върху една карта са четири: \textbf{Open \& edit} отваря изчислението в
калкулатора, \textbf{Export PDF} изтегля го като документ, \textbf{Request support}
отваря формата за запитване с вече прикачен проект, а \textbf{Delete} го изтрива.
Изтриването е окончателно и премахва и запитванията по този проект.

Лентата над списъка съдържа търсене по име, етикет и потребителско свойство, филтър по
етикет и подреждане по последна промяна, по дата на създаване или по име. Бутонът
\textbf{Select} превключва в режим на избор, в който няколко проекта се маркират
едновременно и се изтеглят заедно с \textbf{Export ZIP} или се изтриват с
\textbf{Delete}.

\subsection{Запитване към инженерния отдел}
\label{subsec:ug-support}

Разделът \textbf{Engineering Support} служи за въпрос към инженерния отдел по конкретно
изчисление \figref{fig:ss-support}. Запитване без избран проект не се изпраща: смисълът
на формата е въпросът да носи изчислителния контекст, за да не се налага той да се
преписва и да не се допуска разминаване.

\screenshot{support}{Страница „Engineering Support“: формата за запитване вляво, прикаченото изчисление вдясно}

Формата е в две части. \textbf{Contact details} съдържа име, фамилия, електронна поща,
държава и по желание телефон; името и адресът се попълват от профила. \textbf{Project
and inquiry} съдържа избор на запазен проект, тема (\emph{Result review},
\emph{Supporting structure}, \emph{Attachment method}, \emph{Special structural case},
\emph{Other}) и самия въпрос, до 2800 знака. Съгласието за обработка на изпратените
данни е задължително; съгласието за получаване на рекламни съобщения не е.

Вдясно стои панелът \textbf{Attached calculation}, който показва какво точно ще получи
инженерният отдел: тип съоръжение, височина, схема на закрепване, разстояние между
колоните, наклон, мерна система и водещия товар. Бутонът \textbf{Send inquiry} изпраща
запитването и извежда потвърждение с номер за справка.

\subsection{Техническа документация}
\label{subsec:ug-docs}

Разделът \textbf{Technical Documentation} събира два документа
\figref{fig:ss-docs}. \textbf{Application Manual} е приложното ръководство към товарните
таблици, разгърнато като страници с прелистване (\emph{Previous}, \emph{Next}), избор на
страница, увеличение (\emph{Fit page}, $+$, $-$, \emph{100\%}) и изтегляне като PDF.
\textbf{Standard Attachment Details} е обзорният лист с типовите детайли за закрепване
към бетонов под, бетонова стена, стоманена колона и зидана стена, с увеличение,
придвижване с влачене и изтегляне като PDF.

\screenshot{docs}{Страница „Technical Documentation“: приложното ръководство с прелистване и увеличение}

Приложното ръководство е неразделна част от товарните таблици. Всяка страница с
резултати съдържа изрично указание за това, както и указанието, че товарите са
предварителни и не са за изпълнение.

\subsection{Съобщения и как се процедира}
\label{subsec:ug-messages}

\begin{table}[htbp]
\centering
\caption{Съобщения на системата и действие на потребителя}
\label{tab:ug-messages}
{\fontsize{10}{12}\selectfont
\begin{tabular}{@{}L{5.0cm}L{7.8cm}@{}}
\toprule
\textbf{Съобщение} & \textbf{Значение и действие} \\
\midrule
Select the calculator inputs to view the results. & Някоя от величините не е потвърдена след привеждане. Натиснете отново засегнатата стойност. \\
No table entry for this combination. & За тази комбинация еталонните таблици нямат ред. Системата не екстраполира; изберете съседна стойност или се обърнете към инженерния отдел. \\
Password must be at least 8 characters & Паролата е твърде къса. \\
An account with this email already exists & Профил с този адрес вече съществува; влезте вместо да се регистрирате. \\
Invalid email or password & Данните за вход не съвпадат. Съобщението е едно и също и при несъществуващ профил. \\
Not authenticated & Сесията е изтекла. Влезте отново; изчислението на екрана се запазва. \\
Database unavailable & Услугата за профили е временно недостъпна. Калкулаторът и документацията продължават да работят; запазването на проекти е спряно. \\
Too many attempts, please try again later & Надвишен е допустимият брой опити за вход. Изчакайте и опитайте отново. \\
\bottomrule
\end{tabular}}
\end{table}

\subsection{Работа от телефон}
\label{subsec:ug-mobile}

Уеб приложението работи от телефон без разлика в поведението: същите избори, същите
резултати, същите проекти. Разликата е само в подредбата \figref{fig:phonestrip}.
Панелът с изборите застава над резултата, а схемата под таблицата; разделите от
заглавната лента се подреждат в две колони; лентата с филтри в списъка с проекти става
решетка от етикет и поле, за да не се откъсне надписът от полето, което назовава.
Проверено е, че при всяка ширина от 320 до 2560 px страницата не се разширява настрани
и цялото съдържание е достижимо (§\ref{subsec:viewportres}).

\phonestrip{Същите три страници при ширина 375 px: съдържанието е същото, подредбата е в една колона}
